% GENERATED FILE. Edit canonical lessons, then run npm run generate:books.
% canonical-source-hash: fnv1a64:c160b688c89450b7
% canonical-lessons: TA-C488-paruman, TA-C488-toppai, TA-C488-meli, TA-C488-kun, TA-C488-pattini

\chapter{Weight and Shape}
\label{ch:ta-weight-and-shape}
\hlchaptermodality{488}

\begin{chapteropening}
\textbf{By the end of this chapter:} \emph{I can talk about bulk, a paunch, getting thin, a stoop and going hungry.}

Complete the last lesson of chapter 488: I can talk about bulk, a paunch, getting thin, a stoop and going hungry.
\end{chapteropening}

\section[\texorpdfstring{\ta{பருமன்} (parumaṉ)}{பருமன் (parumaṉ)}]{\ta{பருமன்} (parumaṉ) — fatness; bulk}

\label{lesson:TA-C488-paruman}



\begin{quote}
Before the new one: say the Tamil for to digest, then the Tamil for to burp.
\end{quote}

\subsection*{You'll want to know: \ta{பருமன்}}
\textbf{\ta{பருமன்}} — \emph{parumaṉ} — \textquotedblleft{}fatness; bulk\textquotedblright{}.

\begin{quote}
\ta{உடல்} \ta{பருமன்} \ta{பல} \ta{நோய்களுக்குக்} \ta{காரணம்}.
\end{quote}

\begin{quote}
\emph{uṭal parumaṉ pala nōykaḷukkuk kāraṇam.}
\end{quote}

\begin{quote}
Obesity is the cause of many illnesses.
\end{quote}

\ta{உடல்} \ta{பருமன்}, obesity.

\begin{cousinweb}
\textbf{\ta{பருமன்}} (\emph{parumaṉ}) is native Dravidian, built on \ta{பரு} \emph{paru}, \textquotedblleft{}to become large, bulky\textquotedblright{} (DEDR 3972).
\end{cousinweb}

\begin{grammarlens}[title={What you've built}]
One more word for describing what you lived and felt.
\end{grammarlens}

\subsection*{Your turn}
\begin{itemize}
\raggedright
  \item \emph{Say it:} \emph{parumaṉ}
  \item \emph{Say it:} \emph{parumaṉ}, once more
  \item \emph{Say it:} \emph{uṭal parumaṉ pala nōykaḷukkuk kāraṇam.} — Obesity is the cause of many illnesses.
  \item \emph{Recall:} say the Tamil for to digest, then the Tamil for to burp, then say \emph{parumaṉ} again
\end{itemize}

\begin{tcolorbox}[breakable,colback=teal!4,colframe=teal!35!black,title={Before you move on}]
What does \ta{பருமன்} mean? (\textquotedblleft{}Fatness; bulk\textquotedblright{}.) Say it once more.
\end{tcolorbox}
\section[\texorpdfstring{\ta{தொப்பை} (toppai)}{தொப்பை (toppai)}]{\ta{தொப்பை} (toppai) — a paunch; a pot belly}

\label{lesson:TA-C488-toppai}



\begin{quote}
Before the new one: say the Tamil for to burp, then the Tamil for fatness.
\end{quote}

\subsection*{You'll want to know: \ta{தொப்பை}}
\textbf{\ta{தொப்பை}} — \emph{toppai} — \textquotedblleft{}a paunch; a pot belly\textquotedblright{}.

\begin{quote}
\ta{உடற்பயிற்சி} \ta{செய்யாததால்} \ta{மாமாவுக்குத்} \ta{தொப்பை} \ta{வந்தது}.
\end{quote}

\begin{quote}
\emph{uṭaṟpayiṟci ceyyātatāl māmāvukkut toppai vantatu.}
\end{quote}

\begin{quote}
Uncle got a paunch because he did not exercise.
\end{quote}

Everyday and a little teasing.

\begin{cousinweb}
The origin of \textbf{\ta{தொப்பை}} (\emph{toppai}) is not settled; we could not find it in DEDR.
\end{cousinweb}

\begin{grammarlens}[title={What you've built}]
One more word for describing what you lived and felt.
\end{grammarlens}

\subsection*{Your turn}
\begin{itemize}
\raggedright
  \item \emph{Say it:} \emph{toppai}
  \item \emph{Say it:} \emph{toppai}, once more
  \item \emph{Say it:} \emph{uṭaṟpayiṟci ceyyātatāl māmāvukkut toppai vantatu.} — Uncle got a paunch because he did not exercise.
  \item \emph{Recall:} say the Tamil for to burp, then the Tamil for fatness, then say \emph{toppai} again
\end{itemize}

\begin{tcolorbox}[breakable,colback=teal!4,colframe=teal!35!black,title={Before you move on}]
What does \ta{தொப்பை} mean? (\textquotedblleft{}A paunch; a pot belly\textquotedblright{}.) Say it once more.
\end{tcolorbox}
\section[\texorpdfstring{\ta{மெலி} (meli)}{மெலி (meli)}]{\ta{மெலி} (meli) — to grow thin}

\label{lesson:TA-C488-meli}



\begin{quote}
Before the new one: say the Tamil for fatness, then the Tamil for a paunch.
\end{quote}

\subsection*{You'll want to know: \ta{மெலி}}
\textbf{\ta{மெலி}} — \emph{meli} — \textquotedblleft{}to grow thin\textquotedblright{}.

\begin{quote}
\ta{நோய்க்குப்} \ta{பிறகு} \ta{அவர்} \ta{மெலிந்துவிட்டார்}.
\end{quote}

\begin{quote}
\emph{nōykkup piṟaku avar melintuviṭṭār.}
\end{quote}

\begin{quote}
He has grown thin after his illness.
\end{quote}

Past (I …): \ta{மெலிந்தேன்}, \emph{melintēṉ}.

Past with -\ta{ந்த}-: \ta{மெலிந்தேன்}.

\begin{cousinweb}
\textbf{\ta{மெலி}} (\emph{meli}) is native Dravidian (DEDR 5078), in the family of \ta{மெல்} \emph{mel}, \textquotedblleft{}soft, tender\textquotedblright{}.
\end{cousinweb}

\begin{grammarlens}[title={What you've built}]
One more word for describing what you lived and felt.
\end{grammarlens}

\subsection*{Your turn}
\begin{itemize}
\raggedright
  \item \emph{Say it:} \emph{meli}
  \item \emph{Say it:} \emph{meli}, once more
  \item \emph{Say it:} \emph{nōykkup piṟaku avar melintuviṭṭār.} — He has grown thin after his illness.
  \item \emph{Recall:} say the Tamil for fatness, then the Tamil for a paunch, then say \emph{meli} again
\end{itemize}

\begin{tcolorbox}[breakable,colback=teal!4,colframe=teal!35!black,title={Before you move on}]
What does \ta{மெலி} mean? (\textquotedblleft{}To grow thin\textquotedblright{}.) Say it once more.
\end{tcolorbox}
\section[\texorpdfstring{\ta{கூன்} (kūṉ)}{கூன் (kūṉ)}]{\ta{கூன்} (kūṉ) — a hunch; a stoop}

\label{lesson:TA-C488-kun}



\begin{quote}
Before the new one: say the Tamil for a paunch, then the Tamil for to grow thin.
\end{quote}

\subsection*{You'll want to know: \ta{கூன்}}
\textbf{\ta{கூன்}} — \emph{kūṉ} — \textquotedblleft{}a hunch; a stoop\textquotedblright{}.

\begin{quote}
\ta{வயதானதால்} \ta{பாட்டிக்குக்} \ta{கூன்} \ta{விழுந்தது}.
\end{quote}

\begin{quote}
\emph{vayatāṉatāl pāṭṭikkuk kūṉ viḻuntatu.}
\end{quote}

\begin{quote}
With age grandmother developed a stoop.
\end{quote}

\ta{கூன்} \ta{விழு}, to become hunched.

\begin{cousinweb}
\textbf{\ta{கூன்}} (\emph{kūṉ}) is native Dravidian (DEDR 1927), a bend or a hump.
\end{cousinweb}

\begin{grammarlens}[title={What you've built}]
One more word for describing what you lived and felt.
\end{grammarlens}

\subsection*{Your turn}
\begin{itemize}
\raggedright
  \item \emph{Say it:} \emph{kūṉ}
  \item \emph{Say it:} \emph{kūṉ}, once more
  \item \emph{Say it:} \emph{vayatāṉatāl pāṭṭikkuk kūṉ viḻuntatu.} — With age grandmother developed a stoop.
  \item \emph{Recall:} say the Tamil for a paunch, then the Tamil for to grow thin, then say \emph{kūṉ} again
\end{itemize}

\begin{tcolorbox}[breakable,colback=teal!4,colframe=teal!35!black,title={Before you move on}]
What does \ta{கூன்} mean? (\textquotedblleft{}A hunch; a stoop\textquotedblright{}.) Say it once more.
\end{tcolorbox}
\section[\texorpdfstring{\ta{பட்டினி} (paṭṭiṉi)}{பட்டினி (paṭṭiṉi)}]{\ta{பட்டினி} (paṭṭiṉi) — starvation; going hungry}

\label{lesson:TA-C488-pattini}



\begin{quote}
Before the new one: say the Tamil for to grow thin, then the Tamil for a hunch.
\end{quote}

\subsection*{You'll want to know: \ta{பட்டினி}}
\textbf{\ta{பட்டினி}} — \emph{paṭṭiṉi} — \textquotedblleft{}starvation; going hungry\textquotedblright{}.

\begin{quote}
\ta{வறட்சியால்} \ta{பல} \ta{குடும்பங்கள்} \ta{பட்டினி} \ta{கிடந்தன}.
\end{quote}

\begin{quote}
\emph{vaṟaṭciyāl pala kuṭumpaṅkaḷ paṭṭiṉi kiṭantaṉa.}
\end{quote}

\begin{quote}
Because of the drought many families went hungry.
\end{quote}

\ta{பட்டினி} \ta{கிட}, to go without food.

\begin{cousinweb}
\textbf{\ta{பட்டினி}} (\emph{paṭṭiṉi}) is native Dravidian (DEDR 3872), for fasting and starvation alike.
\end{cousinweb}

\begin{grammarlens}[title={What you've built}]
One more word for describing what you lived and felt.
\end{grammarlens}

\subsection*{Your turn}
\begin{itemize}
\raggedright
  \item \emph{Say it:} \emph{paṭṭiṉi}
  \item \emph{Say it:} \emph{paṭṭiṉi}, once more
  \item \emph{Say it:} \emph{vaṟaṭciyāl pala kuṭumpaṅkaḷ paṭṭiṉi kiṭantaṉa.} — Because of the drought many families went hungry.
  \item \emph{Recall:} say the Tamil for to grow thin, then the Tamil for a hunch, then say \emph{paṭṭiṉi} again
\end{itemize}

\begin{tcolorbox}[breakable,colback=teal!4,colframe=teal!35!black,title={Before you move on}]
What does \ta{பட்டினி} mean? (\textquotedblleft{}Starvation; going hungry\textquotedblright{}.) Say it once more.
\end{tcolorbox}
